\section{Neuro-Symbolic RAG Router}
\label{sec:rag_router}

The bridge between unstructured natural language queries and the strict symbolic solver is our \textbf{Neuro-Symbolic RAG Router}~\cite{edge2024graphrag}. The router performs two tasks: first, it parses the query into logical facts and goals; second, it maps those parsed facts to unique database entities in Neo4j and Qdrant.

\subsection{System Architecture Flowchart}
The complete data flow of Omni-IPS is visualised in the TikZ system architecture flowchart in Figure~\ref{fig:architecture}. It depicts the transformation of user input into structured propositions, database entity mapping, reasoning execution, and streamed explanations.

\begin{figure}[p]
    \centering
    \begin{tikzpicture}[node distance=1.8cm, auto]
        % Nodes
        \node (input) [process, fill=red!10, draw=red!70!black] {\textbf{User Query} \\ (Natural Language)};
        \node (llm_parser) [process, below of=input] {\textbf{LLM Parser Agent} \\ (Extracts Facts \& Goal)};
        \node (bp_scanner) [process, below of=llm_parser] {\textbf{Balanced Paren Scanner} \\ (Commutative Canonicalization)};
        \node (router) [process, below of=bp_scanner, fill=yellow!10, draw=yellow!70!black] {\textbf{Neuro-Symbolic Router}};
        
        \node (qdrant) [database, right=2cm of router] {\textbf{Qdrant Vector DB} \\ (Semantic Index)};
        \node (neo4j) [database, left=2cm of router] {\textbf{Neo4j Graph DB} \\ (Logical Schema)};
        
        \node (solver) [agent, below of=router, minimum width=4cm, minimum height=1.2cm] {\textbf{Symbolic Reasoner} \\ (Forward/Backward Chaining)};
        \node (explain) [process, below of=solver] {\textbf{Explainability Agent} \\ (LLM Proof Generation)};
        \node (ui) [process, below of=explain, fill=green!10, draw=green!70!black] {\textbf{Streamlit UI} \\ (Streaming Output)};

        % Paths
        \draw [arrow] (input) -- (llm_parser);
        \draw [arrow] (llm_parser) -- (bp_scanner);
        \draw [arrow] (bp_scanner) -- (router);
        
        \draw [arrow] (router.east) -- node [above] {Vector Query} (qdrant.west);
        \draw [arrow] (qdrant.west) -- node [below, align=left] {Score $\ge 0.85$ \\ (Fact Matching)} (router.east);
        
        \draw [arrow] (router.west) -- node [above] {Cypher Load} (neo4j.east);
        \draw [arrow] (neo4j.east) -- node [below] {Active Rules} (router.west);
        
        \draw [arrow] (router) -- node [right] {Working Memory} (solver);
        \draw [arrow] (solver) -- node [right] {Execution Trace} (explain);
        \draw [arrow] (explain) -- node [right] {Token Stream} (ui);
    \end{tikzpicture}
    \caption{Omni-IPS System Architecture showing the pipeline from natural language query parsing, dual-database routing (Qdrant & Neo4j), deterministic logical solver reasoning, to streamed markdown explanation generation.}
    \label{fig:architecture}
\end{figure}

\subsection{Nested Balanced Parenthesis Parsing}
Simple regular expression split strategies fail when parsing high-school geometry predicates due to nested parentheses. For instance, in:
\begin{center}
    \texttt{RightAngle(Angle(BAC))}
\end{center}
a naive split on commas and parentheses truncates the trailing parenthesis, resulting in the malformed argument \texttt{Angle(BAC} instead of \texttt{Angle(BAC)}. 

To solve this, we implemented a custom \textbf{Balanced Parenthesis Scanner} that parses variables and predicates in linear $O(N)$ time. The scanner processes the inner text of a relation character by character, tracking a nesting level count $l$. Commas are only treated as parameter delimiters when $l = 0$. The scanner mechanism is formally described in Algorithm~\ref{alg:bp_scanner}.

\begin{algorithm}[h]
\caption{Nested Balanced Parenthesis Parser}
\label{alg:bp_scanner}
\begin{algorithmic}[1]
\Procedure{ParseFactString}{$S$}
    \State $S \gets \text{StripSpaces}(S)$
    \State $idx \gets \text{FindFirst}(S, \text{'('})$
    \If{$idx = -1$ \textbf{or} $\text{LastChar}(S) \neq \text{')'}$}
        \State \textbf{return} $\text{Fact}(\text{relation}=\text{'Atom'}, \text{args}=[S], \text{canonical\_val}=S)$
    \EndIf
    \State $relation \gets S[0 \dots idx-1]$
    \State $inner \gets S[idx+1 \dots \text{Length}(S)-2]$
    \State $args \gets []$, $current\_arg \gets []$, $l \gets 0$
    \For{$char$ \textbf{in} $inner$}
        \If{$char = \text{','}$ \textbf{and} $l = 0$}
            \State $\text{Append}(args, \text{Join}(current\_arg))$
            \State $current\_arg \gets []$
        \Else
            \If{$char = \text{'('}$} \State $l \gets l + 1$
            \Statex \hspace{3.4cm} \textbf{else if} $char = \text{')'}$\textbf{ then} $l \gets l - 1$
            \EndIf
            \State $\text{Append}(current\_arg, char)$
        \EndIf
    \EndFor
    \If{$current\_arg \neq \emptyset$}
        \State $\text{Append}(args, \text{Join}(current\_arg))$
    \EndIf
    \State \textbf{return} $\text{Fact}(relation, args)$
\EndProcedure
\end{algorithmic}
\end{algorithm}

\subsection{Exact and Semantic Routing Pipeline}
Once propositions are extracted, they must be mapped to valid entity IDs in the database. As shown in Figure~\ref{fig:qdrant_db} (sourced from \texttt{assets/qdrant\_vector\_db.png}), Qdrant collections index facts and mathematical rules as vectors. The Neuro-Symbolic Router queries these collections using a two-stage routing pipeline:

\begin{figure}[h]
    \centering
    \includegraphics[width=0.85\textwidth]{qdrant_vector_db.png}
    \caption{Representation of vector collections stored in Qdrant representing the mathematical rules and factual annotations indexed as high-dimensional points.}
    \label{fig:qdrant_db}
\end{figure}

\begin{enumerate}
    \item \textbf{Stage 1: Exact Payload Matching}:
        The router first scrolls the Qdrant database, looking for exact text matches on the payload value property:
        \begin{equation}
            \text{exact\_match}(x) = \{ f \in \text{Qdrant} \mid f.\text{value} = x \}
        \end{equation}
        If a match is found, it is mapped immediately with 100\% confidence, bypassing vector computation.
    \item \textbf{Stage 2: Semantic Dense Vector Search}:
        If an exact match fails, the query string is embedded into a dense vector $q \in \R^{384}$ using a SentenceTransformer. We compute the Cosine similarity against all database facts $v_i \in \R^{384}$:
        \begin{equation}
            \simscore(q, v_i) = \frac{q \cdot v_i}{\|q\| \|v_i\|}
        \end{equation}
        Let $v^*$ be the nearest neighbor: $v^* = \arg\max_{v_i} \simscore(q, v_i)$. 
        The mapping is accepted if and only if it exceeds our high-confidence threshold $\tau_{\text{confidence}}$:
        \begin{equation}
            f(q) = \begin{cases} 
                v^* & \text{if } \simscore(q, v^*) \ge 0.85 \\
                \text{ad-hoc\_fallback}(q) & \text{otherwise}
            \end{cases}
        \end{equation}
    \item \textbf{Stage 3: Ad-Hoc Fallback}:
        If the similarity score falls below $0.85$, the system assumes the user has declared a custom parameter (such as a specific length or unique naming). Instead of returning an error, the router dynamically allocates a new fact ID, e.g., \texttt{geo\_fact\_Parallel(AB, CD)}, and injects it into working memory, ensuring robust support for novel variables.
\end{enumerate}
